半径为 2 的圆 \(\gamma(t)=(2\cos t,2\sin t)\) 的曲率是 \(\kappa=1/2=0.5\)；半径为 1 的圆曲率是 1。曲率是半径的倒数——弯得越急，数字越大。从平面曲线开始，这个量逐步推广到空间曲线（加挠率）、曲面（加第二基本形式）、再到 Gauss 曲率。

上一单元把流形、切空间与黎曼度量建成了一般语言。本单元补上曲率，但只取本书需要的部分——曲线与曲面的经典微分几何。一维和二维的流形在这里完全具体化，所有抽象概念都能算出显式答案；Fisher 度量所在统计流形的弯曲、数据流形的局部形状、以及``流形假设''中``弯曲''一词的确切含义，都需要曲率语言。

学习路径沿维数递进：先看平面曲线怎样用弧长参数化消除参数化的任意性，曲率怎样作为``方向变化率''出现；再看空间曲线多出挠率，Frenet 标架怎样把局部几何装进一个活动坐标系；然后进入曲面，第一基本形式给出内蕴度量，第二基本形式刻画在环境空间中的弯曲方式；最后到 Gauss 曲率与 Theorema Egregium——内蕴几何不依赖嵌入。

读完应能：计算一条曲线的曲率与一个曲面的 Gauss 曲率；说清``圆柱面与平面局部等距''为什么不矛盾；判断``这个流形是弯的''究竟指内蕴曲率还是外在形状。

\subsection{一个具体算例}

悬链线 \(y=\cosh x\) 的曲率：\(\kappa(x)=1/\cosh^2 x\)。在顶点 \(x=0\) 处曲率最大（1），向两侧递减。弧长参数化后 \(|\gamma'(s)|=1\) 恒成立，曲率直接是 \(|\gamma''(s)|\)。

球面 \(S^2_R\) 的 Gauss 曲率 \(K=1/R^2\) 处处为正，Gauss--Bonnet 定理给出 \(\int_{S^2}K\,dA=4\pi\)——曲率的总积分只依赖拓扑（Euler 示性数），不依赖具体度量。平面和圆柱面的 \(K=0\)，总积分也是 0。

\subsection{怎么读这一章}

先从曲线的速度、弧长和转弯程度建立一维直觉，再用第一、第二基本形式区分曲面内部测量与外部弯曲。最后看 Gauss 曲率为何能由内部度量决定，以及这些概念怎样迁移到统计与数据流形。

\subsection{本章篇目}

以下篇目按概念依赖排列；若从具体问题进入，也可以先打开最接近当前困惑的一篇，再沿篇内链接回补。

\begin{itemize}
\tightlist
\item \hyperref[sec:12-Real-Analysis-Support/11-Curve-and-Surface-Geometry/01]{``曲线：弧长参数化与曲率''}；
\item \hyperref[sec:12-Real-Analysis-Support/11-Curve-and-Surface-Geometry/02]{``曲面：第一与第二基本形式''}；
\item \hyperref[sec:12-Real-Analysis-Support/11-Curve-and-Surface-Geometry/03]{``从 Gauss 曲率到统计与数据流形''}。
\end{itemize}
